\documentclass[preview]{standalone}

\usepackage{amsmath}
\usepackage{amssymb}
\usepackage{stellar}
\usepackage{definitions}

\begin{document}

\id{psicologia-apprendimento}
\genpage

\section{Apprendimento}

\begin{snippetdefinition}{apprendimento-definition}{Apprendimento}
    L'\emph{apprendimento} è un processo continuo, basato sull'esperienza, che
    produce un cambiamento relativamente stabile e duraturo nel comportamento,
    nel comportamento potenziale o negli stati rappresentativi di un organismo.
\end{snippetdefinition}

\begin{snippet}{apprendimento-esperienza}
    L'apprendimento può avvenire solo tramite l'esperienza e non coincide con i
    cambiamenti dovuti alla sola maturazione biologica. L'esperienza comprende
    una fase di raccolta, valutazione e rielaborazione delle informazioni e una
    fase di azione volta a produrre risposte capaci di modificare l'ambiente.
\end{snippet}

\includesnpt[width=62\%,src=/snippet/static-psychology/learning-information-memory-action.jpg]{centered-img}

\begin{snippetdefinition}{apprendimento-prestazione-distinction}{Distinzione tra apprendimento e prestazione}
    La \emph{distinzione tra apprendimento e prestazione} riguarda la differenza
    tra ciò che è stato appreso, cioè le potenzialità disponibili, e ciò che
    viene effettivamente messo in atto tramite il comportamento osservabile.
\end{snippetdefinition}

\begin{snippet}{proprieta-apprendimento}
    L'apprendimento è un processo:
    \begin{itemize}
        \item basato su pratica ed esperienza relativamente stabili;
        \item continuo, perché si apprende in modo deliberato o inconsapevole;
        \item collegato sia alle conoscenze procedurali sia alle conoscenze
        dichiarative;
        \item osservabile nelle prestazioni, anche se la prestazione non mostra
        sempre tutto ciò che è stato appreso.
    \end{itemize}
\end{snippet}

\begin{snippetexample}{abituazione-example}{Abituazione}
    Immaginiamo di osservare una fotografia piacevole. Alla prima osservazione
    potremmo provare un'emozione abbastanza intensa; se però osserviamo la
    stessa immagine molte volte in un breve intervallo di tempo, la risposta
    emotiva tende a diventare progressivamente più debole.
\end{snippetexample}

\begin{snippetdefinition}{abituazione-definition}{Abituazione}
    L'\emph{abituazione} è il processo per cui, quando uno stimolo viene
    presentato ripetutamente, si osserva una diminuzione della risposta
    comportamentale.
\end{snippetdefinition}

\begin{snippet}{abituazione-funzione}
    L'abituazione è adattiva perché permette di concentrare le risorse sugli
    eventi nuovi presenti nell'ambiente e riduce lo sforzo comportamentale
    impiegato per rispondere a stimoli già incontrati.
\end{snippet}

\begin{snippetdefinition}{sensibilizzazione-definition}{Sensibilizzazione}
    La \emph{sensibilizzazione} è il processo per cui la risposta a uno stimolo
    presentato ripetutamente diventa più forte invece che più debole.
\end{snippetdefinition}

\begin{snippetexample}{sensibilizzazione-example}{Sensibilizzazione}
    Se una persona viene sottoposta allo stesso stimolo doloroso molte volte in
    un breve intervallo di tempo, può provare più dolore durante l'ultima
    somministrazione rispetto alla prima, anche se l'intensità fisica dello
    stimolo rimane costante.
\end{snippetexample}

\begin{snippet}{abituazione-sensibilizzazione}
    La sensibilizzazione è più probabile quando gli stimoli sono intensi o
    fastidiosi; negli altri casi è più frequente l'abituazione. Entrambi i
    processi mostrano come l'esperienza ripetuta possa modificare la risposta
    dell'organismo.
\end{snippet}

\end{document}
